\documentclass[reqno,11pt]{amsart}
\usepackage{math327-lecture}
\usepackage{needspace}
\lectureinfo{14}{Oct. 7, 2026}{Cosets, Normal Subgroups, and Units Modulo \texorpdfstring{$n$}{n}}

% Left/right cosets and normality: handwritten Lecture 10, Sep. 24, 2025,
% pp. 2--5; restored from the earlier Lecture 13 transcription.
% Units and correspondence: material deferred from Lecture 13.
\begin{document}
\makelecturetitle

\section{Left and right cosets}

Let $H\leq G$. Recall that the left and right cosets represented by
$g\in G$ are
\[
  gH=\{gh:h\in H\},\qquad Hg=\{hg:h\in H\}.
\]
So far we have mostly worked with left cosets. We now compare the two
kinds of cosets.

In general, left cosets are \emph{not} right cosets.

\begin{example}[Left and right cosets in $S_3$]
Write $S_3=\{1,x,x^2,y,xy,x^2y\}$ with
$x^3=y^2=1$ and $yx=x^2y$, and let $H=\{1,y\}$.
Then
\[
\begin{array}{r@{\;=\;}l@{\qquad\qquad}r@{\;=\;}l}
  H & \{1,y\}, & H\cdot1 & \{1,y\},\\[2pt]
  xH & \{x,xy\}, & Hx & \{x,yx\}=\{x,x^2y\},\\[2pt]
  x^2H & \{x^2,x^2y\}, & Hx^2 & \{x^2,yx^2\}=\{x^2,xy\}.
\end{array}
\]
Here $yx^2=(yx)x=x^2yx=x^2\cdot x^2y=xy$. So $xH\neq Hx$ and
$x^2H\neq Hx^2$; in fact neither $xH$ nor $x^2H$ is a right coset of $H$
at all.
\end{example}

However, every statement you can make about left cosets has a right coset
analogue, with the same proof after multiplying on the other side. For
example:
\begin{enumerate}[label=\textbf{(\arabic*)}]
  \item Left cosets partition $G$; so do right cosets.
  \item $g_1H=g_2H$ if and only if $g_2^{-1}g_1\in H$; similarly
  $Hg_1=Hg_2$ if and only if $g_1g_2^{-1}\in H$.
  \item $|gH|=|H|$, via $h\mapsto gh$; similarly $|Hg|=|H|$, via
  $h\mapsto hg$.
\end{enumerate}

What about the index? Is the number of left cosets the same as the number
of right cosets?

\begin{proposition}\label{prop:left-right-bijection}
Let $H\leq G$. The map
\[
  \{\text{left cosets of }H\}\longrightarrow\{\text{right cosets of }H\},
  \qquad gH\longmapsto Hg^{-1},
\]
is a well-defined bijection.
\end{proposition}
\begin{proof}
For a subset $S\subseteq G$ write $S^{-1}=\{s^{-1}:s\in S\}$. Since
$(gh)^{-1}=h^{-1}g^{-1}$ and $H^{-1}=H$,
\[
  (gH)^{-1}=H^{-1}g^{-1}=Hg^{-1}.
\]
So the map is simply $S\mapsto S^{-1}$, which depends only on the set
$gH$ and not on the choice of $g$; hence it is well defined. Applying the
same operation to right cosets, $(Hg)^{-1}=g^{-1}H$, gives the inverse
map. Thus it is a bijection.
\end{proof}

\begin{corollary}
The index $[G:H]$ equals both the number of left cosets and the number of
right cosets of $H$ in $G$. In particular, for finite $G$,
$|G|=[G:H]\,|H|$ however we choose to count.
\end{corollary}

\begin{warningbox}
The natural-looking map $gH\mapsto Hg$ is \emph{not} well defined in
general. In the example above, $xH=xyH$, but
$Hx=\{x,x^2y\}\neq\{xy,x^2\}=Hxy$. (Check: $Hxy=\{xy,yxy\}$ and
$yxy=x^2y^2=x^2$.) This is why we use $g^{-1}$.
\end{warningbox}

\section{Normal subgroups and cosets}

Recall that a subgroup $H\leq G$ is \emph{normal}, written
$H\trianglelefteq G$, if $gHg^{-1}\subseteq H$ for all $g\in G$, or
equivalently $gHg^{-1}=H$ for all $g\in G$.

Although left cosets are not right cosets in general, they are for normal
subgroups.

\begin{lemma}\label{lem:normal-left-right}
If $H\trianglelefteq G$, then every left coset of $H$ is also a right
coset. In fact, $gH=Hg$ for all $g\in G$.
\end{lemma}
\begin{proof}
Multiplying $gHg^{-1}=H$ on the right by $g$ gives $gH=Hg$.
\end{proof}

\begin{lemma}\label{lem:which-right-coset}
If a left coset $gH$ equals some right coset $Hg'$, then in fact
$Hg'=Hg$.
\end{lemma}
\begin{proof}
Both the left cosets and the right cosets of $H$ partition $G$. The left
coset $gH$ and the right coset $Hg$ both contain $g=ge_G=e_Gg$, so
$Hg'=gH$ meets $Hg$. Two right cosets that meet are equal, so $Hg'=Hg$.
\end{proof}

\begin{corollary}
If every left coset of $H$ is also some right coset, then $H$ is normal
in $G$.
\end{corollary}
\begin{proof}
Let $g\in G$. By hypothesis $gH$ is a right coset, so by
Lemma~\ref{lem:which-right-coset}, $gH=Hg$. Multiplying on the right by
$g^{-1}$ gives $gHg^{-1}=H$. This holds for all $g\in G$.
\end{proof}

Combining the last two results:
\begin{align*}
  H\trianglelefteq G
  &\quad\Longleftrightarrow\quad gH=Hg\text{ for all }g\in G\\
  &\quad\Longleftrightarrow\quad
  \text{every left coset of $H$ is a right coset.}
\end{align*}
For example, $H=\{1,y\}$ is not normal in $S_3$, matching the computation
$xH\neq Hx$ above.

\section{Example: units modulo \texorpdfstring{$n$}{n} and Fermat's little theorem}

Let $n\geq2$. The set $\Z/n\Z$ of residue classes modulo $n$ is a group
under addition, with identity $\overline0$. Recall that
\[
  a\equiv b\pmod n\quad\text{means}\quad n\mid(b-a).
\]
Now consider the operation of \emph{multiplication} on this set,
$\overline a\cdot\overline b=\overline{ab}$. It is not a group under
multiplication (for instance $\overline0$ has no inverse), but we can
ask: for which $a$ does the congruence
\begin{equation}\label{eq:unit}
  ax\equiv1\pmod n
\end{equation}
have a solution $x\in\Z/n\Z$?

A solution means $ax-1=bn$ for some $b\in\Z$, i.e.
\[
  ax+(-b)n=1.
\]
Such an equation has an integer solution if and only if
$\gcd(a,n)=1$: if $d$ divides $a$ and $n$ then it divides $1$;
conversely, $\Z a+\Z n=\Z\gcd(a,n)$, so if $\gcd(a,n)=1$ then $1$ is an
integer combination of $a$ and $n$. Note that $\gcd(a,n)$ depends only on
$\overline a$, since $\gcd(a+kn,n)=\gcd(a,n)$.

\begin{define}[Units modulo $n$]
The \emph{group of units modulo $n$} is
\[
  (\Z/n\Z)^\times
  =\{\overline a\in\Z/n\Z:\ ax\equiv1\ (\mathrm{mod}\ n)\text{ has a
  solution}\}
  =\{\overline a\in\Z/n\Z:\gcd(a,n)=1\},
\]
with the operation of multiplication.
\end{define}

This is not just a set but a group. The key point is closure: if
$\overline a,\overline b\in(\Z/n\Z)^\times$, choose solutions
$ax_1\equiv1$ and $bx_2\equiv1\pmod n$. Then $y=x_2x_1$ satisfies
\[
  (ab)y=a(bx_2)x_1\equiv ax_1\equiv1\pmod n,
\]
so $\overline{ab}\in(\Z/n\Z)^\times$. The remaining axioms are checked
similarly.

\begin{exercise}
Complete the verification that $(\Z/n\Z)^\times$ is an abelian group:
check associativity, that $\overline1$ is the identity, and that the
solution $\overline x$ of \eqref{eq:unit} is an inverse of $\overline a$
lying in $(\Z/n\Z)^\times$.
\end{exercise}

The order of this group is
\[
  |(\Z/n\Z)^\times|=\varphi(n)=\#\{1\leq a\leq n:\gcd(a,n)=1\},
\]
\emph{Euler's $\varphi$-function}. For example,
$(\Z/8\Z)^\times=\{\overline1,\overline3,\overline5,\overline7\}$ and
$\varphi(8)=4$.

If $n=p$ is prime, every $a$ with $1\leq a\leq p-1$ is coprime to $p$, so
$\varphi(p)=p-1$ and $(\Z/p\Z)^\times$ has order $p-1$. Lagrange's theorem
now gives a classical result from number theory.

\begin{theorem}[Fermat's little theorem]
Let $p$ be prime and $a\in\Z$ with $p\nmid a$. Then
\[
  a^{p-1}\equiv1\pmod p.
\]
\end{theorem}
\begin{proof}
Since $p\nmid a$, $\overline a\in(\Z/p\Z)^\times$, a group of order
$p-1$. By Lagrange's theorem the order $m$ of $\overline a$ divides
$p-1$, say $p-1=mk$. Then $\overline a^{\,p-1}=(\overline a^{\,m})^k
=\overline1$.
\end{proof}

The same argument applied to $(\Z/n\Z)^\times$ gives \emph{Euler's
theorem}: if $\gcd(a,n)=1$, then $a^{\varphi(n)}\equiv1\pmod n$.

\begin{remark}
The group $(\Z/p\Z)^\times$ is in fact cyclic. A generator is called a
\emph{primitive root} modulo $p$. For example, modulo $7$ the powers of
$3$ are
\[
  3,\ 2,\ 6,\ 4,\ 5,\ 1,
\]
so $\overline3$ generates $(\Z/7\Z)^\times$. We will not prove this fact
in general.
\end{remark}

\section{The correspondence theorem}

Let $\varphi:G\to\mathcal G$ be a \emph{surjective} homomorphism of groups,
and let $K=\ker(\varphi)\trianglelefteq G$.

\begin{theorem}[Correspondence theorem]\label{thm:correspondence}
\leavevmode
\begin{enumerate}[label=\textbf{(\arabic*)}]
  \item There is a bijective correspondence
  \[
    \{\text{subgroups }H\leq G\text{ with }K\leq H\}
    \quad\longleftrightarrow\quad
    \{\text{subgroups }\mathcal H\leq\mathcal G\},
  \]
  given by
  \[
    H\longmapsto\varphi(H)=\mathcal H,
    \qquad
    \mathcal H\longmapsto\varphi^{-1}(\mathcal H).
  \]
  \item Under this correspondence,
  $H\trianglelefteq G$ if and only if
  $\mathcal H\trianglelefteq\mathcal G$.
  \item If $H$ is finite, then $|H|=|\mathcal H|\cdot|K|$.
\end{enumerate}
\end{theorem}

Here $\varphi^{-1}(\mathcal H)=\{g\in G:\varphi(g)\in\mathcal H\}$ is the
preimage. Part~(3) follows from what we already know: the restriction
$\varphi|_H:H\to\mathcal H$ is a surjective homomorphism with kernel
$H\cap K=K$, so the counting formula for homomorphisms from Lecture~13
gives $|H|=|K|\cdot|\mathcal H|$.

\end{document}
